\chapter*{Résumé}
\addstarredchapter{Résumé}

\initial{L}a conversion du trouble cognitif léger (MCI) vers la démence est une question clinique et de santé publique centrale. Les modèles d'apprentissage automatique multimodaux affichent souvent une AUC supérieure à 0,8, mais environ deux tiers sont entraînés sur ADNI et seulement 10~\% sont validés en externe ; beaucoup souffrent de fuites longitudinales et de définitions de phénotype instables. Ce stage de deux mois (24 juillet -- 24 septembre 2026), mené à l'UMMISCO en collaboration avec le Glenn Biggs Institute, vise un modèle de risque de conversion à 2 et 5 ans, résistant aux fuites, validé hors site, et un \emph{paquet de données minimal} encore utile.

Le travail livré n'est pas un AUC interne de plus. Il consiste en : (i) une cartographie des cohortes NACC, ADNI et Framingham et de leurs voies d'accès ; (ii) une revue de littérature (CARD/NIH, MOIRA, Kolachalama, leakage, PROBAST) ; (iii) un plan d'analyse au format du laboratoire Fongang (trois aims, phénotypes à figer, modèles emboîtés 1--4, calibration, DCA, SHAP, équité). L'hypothèse de travail est que les extraits sont fournis par les encadrants sous DUA. L'unité d'analyse est le participant à une visite index MCI ; toutes ses visites restent dans la même partition. Les profils incomplets (sans TEP ni LCR) sont conservés.

\paragraph{Mots-clés} : MCI, conversion vers la démence, apprentissage multimodal, fuite de données, validation externe, NACC, ADNI, Framingham, paquet de données minimal.
